% Generated draft: requires mathematical and source review.
\documentclass[11pt]{article}
\usepackage[margin=25mm]{geometry}
\usepackage{fontspec}
\setmainfont{Latin Modern Roman}
\usepackage{amsmath,amssymb,mathtools,unicode-math}
\setmathfont{Latin Modern Math}
\usepackage{enumitem,needspace,xurl,hyperref}
\setlength{\parindent}{0pt}
\setlength{\parskip}{5pt}
\setlength{\emergencystretch}{3em}
\newcommand{\N}{\mathbb{N}}
\newcommand{\Z}{\mathbb{Z}}
\newcommand{\Q}{\mathbb{Q}}
\newcommand{\R}{\mathbb{R}}
\newcommand{\C}{\mathbb{C}}
\newcommand{\F}{\mathbb{F}}
\newcommand{\A}{\mathbb{A}}
\newcommand{\PP}{\mathbb{P}}
\newcommand{\E}{\mathbb{E}}
\newcommand{\eps}{\varepsilon}
\newcommand{\dd}{\,\mathrm{d}}
\begin{document}
\section*{UnsolvedMath ID: 30002930}
\textbf{Source problem code:} catalogue entry
\subsection{Definitions and mathematical statement}
\textbf{Source-level question and unresolved formulation.} The supplied record asks for moment bounds for a stochastic corrector and a corresponding quantitative two-scale expansion in stochastic homogenization. It does not specify a stochastic model or enough hypotheses to determine a single quantified mathematical conjecture. The notation below describes the requested estimates without selecting missing hypotheses.

Let \((\Omega,\mathcal F,\mathbb P)\) denote the probability space of a model to be specified, let \(D\) denote its spatial domain, and let \(\phi(x,\omega)\) denote a chosen, normalized stochastic corrector. These objects, including the value space and its norm, must be supplied by the intended model. For a moment order \(p>0\), define the extended-valued moment
\[
M_p(x):=\mathbb E\bigl[\|\phi(x,\cdot)\|^p\bigr].
\]
The first requested conclusion is either a spatially uniform bound
\[
\sup_{x\in D}M_p(x)\le C_p<\infty,
\]
or a precise description of its divergence in an appropriate spatial or regularization limit. The source does not specify the moment orders, the limiting parameter, whether moments of the corrector itself or of its increments are intended, or whether a precise divergence rate means an upper bound, matching upper and lower bounds, or an asymptotic equivalence. Corrector normalization must be fixed before an estimate for its values can be interpreted.

For the second requested conclusion, let \(u_\varepsilon\) denote a microscopic solution and \(u_\varepsilon^{(2)}\) its intended two-scale approximation, once the equation, data, correctors, and any boundary corrections have been specified. A quantitative expansion would require a specified estimate of the schematic form
\[
\|u_\varepsilon-u_\varepsilon^{(2)}\|_{\mathcal X}
\le C\,r(\varepsilon),
\qquad r(\varepsilon)\longrightarrow0
\quad\text{as }\varepsilon\downarrow0,
\]
with the error norm, probabilistic interpretation, rate, and dependence of the constant explicitly determined. This display is a template for the requested conclusion, not an asserted estimate or a choice of solution framework.

\textbf{Research question.} Determine the hypotheses and quantifiers under which the requested corrector moment estimates and corresponding quantitative two-scale expansion hold, and determine the relevant bounds and rates. The supplied dated assessment refers to a general ergodic formulation, but no precise ergodicity assumption or coefficient class is stated. Consequently, the record cannot presently be converted into a unique theorem-style statement. In particular, no universal rate under ergodicity alone is asserted here.
\subsection{Short English statement}
The problem asks how large stochastic correctors can be in moments, whether those moments stay bounded in space or grow at a precise rate, and how such estimates control a two-scale approximation. The model, normalization, moment orders, equations, boundary conditions, error norm, and rates remain unspecified. The supplied partially solved label is dated curation, not an independently verified mathematical conclusion.
\subsection{Sources}
\begin{itemize}
\item Local UnsolvedMath registry, \texttt{data/unsolvedmath/problems.json}, permanent ID \(30002930\), title: \emph{Uniform Corrector Moments and Quantitative Stochastic Homogenization}; source problem code \texttt{catalogue entry}. The supplied record belongs to the Oberwolfach Reports Open Problems collection and lists proposed year \(2015\). No individual proposer is identified.
\item Supplied report attribution: \emph{Interplay of Analysis and Probability in Applied Mathematics (2016)}, Oberwolfach Reports, pp.\ 1989--2064, DOI \texttt{10.4171/owr/2015/35}. \url{https://doi.org/10.4171/owr/2015/35}. The background describes Oberwolfach Reports as documenting mathematical workshops and their open-problem discussions. No problem-specific page or speaker attribution is supplied. The citation's \(2016\) designation and the record's \(2015\) proposed year are retained without resolving their relationship.
\item Supplied literature reference: Gloria--Otto, quantitative corrector estimates. \url{https://arxiv.org/abs/1409.0801}.
\item Supplied literature reference: Gloria--Otto, corrector rates and two-scale expansion. \url{https://arxiv.org/abs/1510.08290}.
\item Literature assessment dated exactly \texttt{2026-08-21}: the supplied status is \texttt{partially\_solved}. The supplied assessment states that optimal corrector moment and growth estimates and quantitative two-scale expansions are available under spectral-gap, finite-range, or related quantitative assumptions, while the general ergodic statement in the record is not fully resolved. This is supplied dated curation derived from the OWR material, not an independently verified literature assessment. The background expressly calls for independent verification of time-sensitive status and cited claims.
\item The registry labels the statement \texttt{exact} and describes it as clear, self-contained, publication-ready, and retained byte-for-byte. These are supplied editorial qualifications; they do not resolve the missing mathematical hypotheses listed in this draft.
\end{itemize}
\end{document}
